\section{Theorem 4}\label{s:thm}

\lem{5}{For every such agent, a life of positive meaning is available.}

\paragraph{The sentence as a whole.} Theorem~4 is the only result of EXI-01-a, and it states the second half of the family conclusion: meaning in the sense of Definition~1 is available to everyone. The first half, that life is meaningful to the extent that one engages with projects of objective worth, is stated by Definition~1 and is not a theorem. The text thus divides the family conclusion between a definition and a theorem, and Remark~5 treats the theorem, with its definition, as answering both of the questions the title raises (\Lm{9}).

The theorem has the logical form $\forall a\,\exists l\,(l \text{ is a life of } a \wedge M(l) > 0)$: for every agent in the domain, there is a life of hers that has positive meaning. It is weaker than it sounds in three ways and stronger in one. It is weaker because \enquote{positive} means no more than \enquote{greater than zero}, because \enquote{available} means no more than \enquote{possible for her from her situation}, and because it says nothing about actual lives. It is stronger because its quantifier is universal and admits no exception.

\lem{5.1}{For every such agent}

\enquote{Such} refers back to the domain of Assumption~3, agents capable of intentional action. The theorem therefore has exactly the scope of the assumption from which it is derived, and no wider. In particular, it has nothing to say about beings that are not agents in this sense; the comment on \Lm{4.2} has discussed what follows for the reading of \enquote{everyone}. The anaphor also shows that Theorem~4 is meant to be read together with Assumption~3, and the proof confirms this: the theorem is, in effect, Assumption~3 transposed from projects to lives.

\lem{5.2}{a life of positive meaning}

\paragraph{\enquote{a life}.} The indefinite article is important. The theorem does not say that the agent's life has positive meaning, or that it will; it says that there is a life, among those available to her, that does. Her actual life may be one of the others.

\paragraph{\enquote{positive meaning}.} The phrase is defined through Definition~1: a life has positive meaning if it is meaningful to some degree greater than zero. The theorem claims nothing about how much. This is the source of the objection, which EXI-01-d considers, that availability so secured is trivial: if a life of the barest positive meaning counts, it is no surprise that one is available to everyone.

The reply has three parts. First, the claim is not trivial, since it is denied by every view on which meaning requires success, talent, a cosmic purpose or an exceptional life, and these views have been widely held. EXI-17-a argues against the first, EXI-01-d against the second, and EXI-11-a against the third (its Theorem~3.3). Second, a minimal claim was what the question required. The question whether meaning is available to everyone is a question whether anyone is shut out of it, and to show that no one is, one needs to show only that a positive degree is available to each. Whether a high degree is available to each is a different question, to which the answer is no. Third, \enquote{positive} meaning in the sense of Definition~1 is not negligible meaning. A life that contains sustained, caring engagement with the well-being of one other person, or with the understanding of one part of the world, has positive meaning, and reflection does not judge such a life to be meaningful only in a degenerate sense. The minimal case is not a borderline case.

\paragraph{Positive meaning and equality.} The theorem implies nothing about equality of meaning. EXI-01-d draws the distinction: availability is universal, degree is not, and it would be an invalid inference from the first to the claim that all lives are equally meaningful, or that deprivation is harmless to meaning (its section on what availability does not imply). The comment on \Lm{4.3} has noted that the minimal form of Assumption~3 is compatible with great inequality in the range of projects open to different agents. Theorem~4 inherits that compatibility. It is a floor, not a level.

\lem{5.3}{is available}

\paragraph{Availability and actuality.} \enquote{Available} is the theorem's modal operator. A life of positive meaning is available to an agent if it is among the lives she can lead. The text does not define availability, and the proof relies on an unstated principle connecting it with Assumption~3's ability to engage (Proposition~\ref{p:bridge}). What the word excludes is clear: it does not say that the agent leads, will lead or is likely to lead such a life. Remark~5 makes the point explicit (\Lm{10}).

\paragraph{Availability at a time.} A question the text does not address is whether availability is indexed to times. Read without an index, the theorem says that among the lives an agent could have led, from the beginning, some have positive meaning. That is weaker than what the family conclusion seems to intend, which is that meaning is available to her now, whatever her situation. The indexed reading says: for every agent and every time at which she is an agent, there is a continuation of her life from that time in which it has positive meaning. On the indexed reading, the theorem says that it is never too late; on the unindexed reading, it says only that it was not always too late.

The proof supports the indexed reading if Assumption~3 is read as indexed: if at every time at which she is an agent there is a project of objective worth in which she is then able to engage. EXI-01-c's argument supports the indexed assumption, since persons and the world as an object of understanding are present at every time in a human life. And Definition~1 permits the inference: a life that, from some time on, includes sustained engagement with a worthwhile project consists to a positive extent in such engagement, on any measure satisfying Proposition~\ref{p:neutral}. We therefore read the theorem as indexed. It says that a life of positive meaning remains available to an agent for as long as she remains one.

\paragraph{Availability and luck.} The indexed reading connects Theorem~4 with EXI-17-a's Corollary~3.6, that meaning is not hostage to luck in outcomes, which EXI-17-a says supports, without by itself establishing, the availability claim of EXI-01. Luck enters a life in two ways: in the circumstances from which one acts, and in the outcomes of one's actions. EXI-17-a addresses the second: since meaning does not require success, bad luck in outcomes cannot remove it. Theorem~4, read with EXI-01-c, addresses the first: since every situation contains objects of worth within reach, bad luck in circumstances cannot remove the possibility of meaning either, though it can lower its attainable degree. Together they give the sense in which meaning, on the account of EXI-01, is immune to fortune: its possibility is, and its degree is not. \citet{Williams1981b} argued that the justification of some choices depends on luck in outcomes, with the example of a painter who abandons his family for his art and whose choice is vindicated only if he succeeds. If meaning depended on such vindication, the painter's life would be meaningful only if he were lucky. On Definition~1 it is meaningful to the extent of his engagement with painting, whether he succeeds or not; what depends on luck is the justification of his choice and the value his work adds to the world, not the meaning of his life. Section~\ref{s:whole} returns to the case.

\subsection{Availability on the rival views}\label{s:availrivals}

Theorem~4 is distinctive of EXI-01-a only if the rival accounts of meaning do not deliver it, or deliver it at a cost Definition~1 avoids. We consider the rival schemata of Section~\ref{s:rivals} in turn.

On subjectivism (R1), meaning is available to every agent trivially, since every agent capable of intentional action wants something and can do something she wants. Availability is secured, but at the cost of the cases: the grass counter and Taylor's Sisyphus have as much meaning as anyone. Subjectivism makes meaning universally available by making it indiscriminate.

On pure objectivism (R2), meaning is available to every agent whose activity can be directed at objective value, which, given the reach of persons and of the world as an object of understanding, is every agent. Availability is secured at the cost of the alienation objection (Objection~\ref{o:alien}): the meaning of a life could be present without its subject's caring about anything, and the subject who discovers that her indifferent activity was meaningful has discovered something about the world rather than about her life.

On the good cause schema (R3), meaning is available only to those whose activity can cause good outcomes, and whether it does depends on luck in outcomes. EXI-17-a's argument applies directly: the unsuccessful mathematician has no meaning in respect of her work on this schema, and whether a given agent's efforts will produce good is not in her power. The schema does not secure availability, and the failure is principled: on it, meaning is hostage to luck in outcomes.

On the transcendence schema (R4), meaning is available to those who can make a difference to something of value beyond themselves. Whether one can make a difference is less dependent on luck than whether one can cause good outcomes, since a difference can be small, but the prisoner in solitary confinement and the listener who produces nothing (Section~\ref{s:hard}; EXI-17-a, Lemma~3.2) make no difference that leaves traces, and the schema does not secure availability for them.

On the purpose schema (R5), meaning is available to every agent if and only if a creator has assigned to every human life a purpose that it can fulfil. EXI-11 holds that meaning does not require a cosmic purpose; on the purpose schema, availability is secured only if theism is true and the purposes assigned are within every agent's reach, and EXI-11-a's Lemma~3.1 shows that even then the assigned purpose confers meaning only if it is good.

\begin{proposition}[Availability without its costs]\label{p:availcost}
Of the six schemata, Definition~1 is the only one that secures the availability of meaning to every agent capable of intentional action without making meaning indiscriminate, without making it independent of the subject's care, without making it depend on luck in outcomes, and without making it depend on theism.
\end{proposition}

\begin{proof}[Argument]
By the discussion above: subjectivism secures availability and makes meaning indiscriminate; pure objectivism secures it and makes meaning independent of care; the good cause and transcendence schemata do not secure it, because they make meaning depend on outcomes or traces; the purpose schema secures it only if theism is true. Definition~1 secures it by Theorem~4, with Assumption~3 supported by EXI-01-c, and it avoids the four costs by its two components: worth excludes indiscriminateness, engagement excludes independence of care, the constitutive reading of \enquote{consists in} excludes dependence on outcomes, and Assumption~2 makes no reference to a creator.
\end{proof}

The proposition gives the strongest reason for the form of Definition~1. Each of its components is required to avoid one of the costs, and the two together secure availability without any of them. The family conclusion's two halves, which the text divides between a definition and a theorem, are in this sense connected: the account of what meaning is was not chosen to make meaning available, but it is the account on which meaning is available for the right reasons.
